\documentclass[10pt,a4paper]{amsart}
\usepackage[margin=19mm]{geometry}
\usepackage{amsmath,amssymb,mathtools}
\usepackage[scr=rsfs,cal=euler]{mathalfa}
\usepackage{xfrac}
\usepackage[unicode,hidelinks,bookmarks=false]{hyperref}
\newcommand{\ie}{i.e.\ }
\newcommand{\cf}{cf.\ }
\newcommand{\resp}{respectively\ }
\newcommand{\Subsection}{\subsection}
\providecommand{\nobreakdash}{\nobreak\hbox{-}}
\setlength{\emergencystretch}{2em}
\multlinegap0pt
\usepackage{bm}
\usepackage{booktabs}
\usepackage{xcolor}
\usepackage{amssymb}
\usepackage{mathtools}
\newtheorem{assumption}{Assumption}
\newtheorem{lemma}{Lemma}
\newcommand{\barg}{\bar{g}}
\newcommand{\bx}{\boldsymbol{x}}
\title{Why is Normalization Preferred? A Worst-Case Complexity Theory for Stochastically Preconditioned SGD under Heavy-Tailed Noise}
\author{Yuchen Fang, James Demmel, Javad Lavaei}
\date{Source: arXiv:2602.13413v1 (CC BY 4.0)}
\begin{document}
\maketitle

\noindent\emph{Source and licence.} Why is Normalization Preferred? A Worst-Case Complexity Theory for Stochastically Preconditioned SGD under Heavy-Tailed Noise. Version arXiv:2602.13413v1 (CC BY 4.0), distributed by arXiv under the Creative Commons Attribution 4.0 International licence, http://creativecommons.org/licenses/by/4.0/, as recorded in the arXiv licence metadata for this exact version and shown on the abstract page https://arxiv.org/abs/2602.13413v1. Full licence notice: \texttt{licenses/arxiv-2602.13413-CC-BY-4.0.txt}.\par\smallskip

\noindent\textit{Editorial scope.} Counted unit: the article's lemma lemma (source0.tex lines 962--971). The article's complete original proof of it is delivered with it (source0.tex lines 974--1139, one \texttt{proof} environment of 166 lines). No mathematics is written, repaired or re-derived: the statement and the proof are the article's own lines, in the article's own notation, and no retained line is substituted or re-flowed.  Retained uncounted as the source-local context the proof needs: lines 416--421; lines 423--430; lines 578--580.  Nothing was omitted from any retained span, so the package carries no omission record.  No retained line contains a citation, so the wrapper carries no bibliography and no bibliography entry.  No retained line contains a figure environment, an image-inclusion command or drawing code, so no image is delivered and the package declares no asset dependency.  Editorial formatting only, in addition: an AMS \texttt{amsart} preamble that restates the article's own declarations and macros used by the retained lines, namely the environments \texttt{assumption}, \texttt{lemma} on separate counters, together with the standard packages those lines need.  The retained environments print in this document as Assumption 1, Lemma 1 in order of appearance, whereas the article prints the same items with the numbering of its own full document.  The complete native source of the article as distributed in the arXiv e-print for this exact version is bundled unchanged as \texttt{sources/194634/source0.tex}.
\begin{assumption}[$L$-smoothness]\label{assumption2}
There exists $L > 0$ such that for any $\bx, \boldsymbol{y} \in \mathbb{R}^d$, we have
\begin{equation*}
  \|\nabla f(\boldsymbol{x}) - \nabla f(\boldsymbol{y})\| \leq L \|\bx - \boldsymbol{y}\|.  
\end{equation*}
\end{assumption}

\begin{assumption}[$p$-BCM condition]\label{assumption3}
At the $k$-th iteration, the gradient estimate $\barg_k$ is unbiased with bounded $p$-th moment, that is,
$\mathbb{E}_k[\barg_k ] = \nabla f(\bx_k)$ and
\begin{equation*}
\mathbb{E}_k \big[ \|\barg_k - \nabla f(\bx_k)\|^p \big] \leq \sigma^p
\end{equation*}
for some $p \in (1,2]$ and $\sigma \geq 0$.
\end{assumption}

\begin{assumption}[Uniform Boundedness of $D_k$]\label{assumption5}
    There exist deterministic constants $m_D,M_D>0$ such that $m_D\cdot I \preceq D_k \preceq M_D\cdot I$ for $\forall k \leq T$.
\end{assumption}

\begin{lemma}
    Under Assumptions \ref{assumption2}, \ref{assumption3}, and \ref{assumption5}, when $\|\nabla f(\bx_k)\|\geq \tau/2$, if we select $\eta<\frac{m_D}{24L }$ and 
    \begin{equation*}
        \tau = \max\left\{2,  4\cdot 3^{1/p}\cdot \sigma\left(\frac{M_D}{m_D}\right)^{\frac{p+1}{2p}}, 64 \sigma \frac{M_D}{m_D} \right\},
    \end{equation*}
    we have
    \begin{align*}
    \mathbb{E}_k[ f(\bx_{k+1}) ] \le f(\bx_k) - \frac{1}{12} \eta m_D \|\nabla f(\bx_k)\|.
\end{align*}
\end{lemma}

\begin{proof}

By the Taylor expansion, we have
    \begin{align*}
    f(\bx_{k+1})  \le f(\bx_k) + \nabla f(\bx_k)^T \Delta \bx_k + \frac{L }{2} \|\Delta \bx_k\|^2  = f(\bx_k) - \eta  \nabla f(\bx_k)^T D_k\widehat{\barg}_k + \frac{\eta^2 L }{2} \|D_k\widehat{\barg}_k\|^2.
\end{align*}
We first consider the conditional expectation of the 
term $\nabla f(\bx_k)^T D_k\widehat{\barg}_k$. We have
\begin{align*}
\mathbb{E}_k\left[\nabla f(x_k)^T D_k \widehat{\barg}_k \right] = \mathbb{E}_k\left[\nabla f(x_k)^T D_k \bar{g}_k \cdot \boldsymbol{1}_{(\|\bar{g}_k\| \le \tau)}\right]  + \tau \cdot \mathbb{E}_k\left[\nabla f(x_k)^T D_k \frac{\bar{g}_k}{\|\bar{g}_k\|} \cdot \boldsymbol{1}_{(\|\bar{g}_k\| > \tau)} \right]
\end{align*}

We focusing on the first term:
\begin{align*}
& \mathbb{E}_k\left[\nabla f(x_k)^T D_k \bar{g}_k \cdot \boldsymbol{1}_{(\|\bar{g}_k\| \le \tau)}\right] \\
&= \mathbb{E}_k\left[\nabla f(x_k)^T D_k (\nabla f(x_k) + \bar{g}_k - \nabla f(x_k)) \cdot \boldsymbol{1}_{(\|\bar{g}_k\| \le \tau)} \right] \\
&= \mathbb{E}_k\left[\nabla f(x_k)^T D_k \nabla f(x_k) + \nabla f(x_k)^T D_k (\bar{g}_k - \nabla f(x_k)) \cdot \boldsymbol{1}_{(\|\bar{g}_k\| \le \tau)} \right]\\
& = \mathbb{E}_k\left[\nabla f(x_k)^T D_k \nabla f(x_k)\right] + \mathbb{E}_k\left[\nabla f(x_k)^T D_k (\bar{g}_k - \nabla f(x_k)) \cdot \boldsymbol{1}_{(\|\bar{g}_k\| \le \tau)}\right]\\
& \geq \mathbb{E}_k\left[\nabla f(x_k)^T D_k \nabla f(x_k)\right] - \mathbb{E}_k\left[\|D_k^{1/2}\nabla f(x_k)\| \|D_k^{1/2} (\bar{g}_k - \nabla f(x_k))\| \cdot \boldsymbol{1}_{(\|\bar{g}_k\| \le \tau)} \right]\\
& = \mathbb{E}_k\left[\nabla f(x_k)^T D_k \nabla f(x_k)\right] \\
& \quad - \mathbb{E}_k\left[\|D_k^{1/2}\nabla f(x_k)\| \|D_k^{1/2} (\bar{g}_k - \nabla f(x_k))\| \cdot \boldsymbol{1}_{(\|\bar{g}_k\| \le \tau,\|D_k^{1/2}(\bar{g}_k - \nabla f(x_k))\|\leq \sqrt{m_D}\cdot\tau/4)}\right]\\
& \quad - \mathbb{E}_k\left[\|D_k^{1/2}\nabla f(x_k)\| \|D_k^{1/2} (\bar{g}_k - \nabla f(x_k))\| \cdot \boldsymbol{1}_{(\|\bar{g}_k\| \le \tau,\|D_k^{1/2}(\bar{g}_k - \nabla f(x_k))\|> \sqrt{m_D}\cdot\tau/4)}\right] .
\end{align*}

We first consider the case when $\|\bar{g}_k\| \le \tau,\|D_k^{1/2}(\bar{g}_k - \nabla f(x_k))\|\leq \sqrt{m_D}\cdot\tau/4$.

Since $\|\nabla f(\bx_k)\|>\tau/2$, we have
\begin{equation*}
    \|D_k^{1/2}\nabla f(\bx_k)\|^2 = \nabla f(\bx_k)^TD_k\nabla f(\bx_k)\geq m_D\|\nabla f(\bx_k)\|^2\geq \frac{\tau^2m_D}{4},
\end{equation*}
which imply
\begin{equation*}
    \|D_k^{1/2}\nabla f(\bx_k)\| \geq \frac{\tau \sqrt{m_D}}{2}, \quad \|D_k^{1/2}(\bar{g}_k - \nabla f(x_k))\|\leq \sqrt{m_D}\cdot\tau/4 \leq \frac{1}{2}\|D_k^{1/2}\nabla f(\bx_k)\|.
\end{equation*}
Therefore we find
\begin{equation*}
    \|D_k^{1/2}\nabla f(x_k)\| \|D_k^{1/2} (\bar{g}_k - \nabla f(x_k))\|\leq \frac{1}{2}\|D_k^{1/2}\nabla f(\bx_k)\|^2.
\end{equation*}
We further have
\begin{align*}
& \mathbb{E}_k\left[\|D_k^{1/2}\nabla f(x_k)\| \|D_k^{1/2} (\bar{g}_k - \nabla f(x_k))\| \cdot \boldsymbol{1}_{(\|\bar{g}_k\| \le \tau,\|D_k^{1/2}(\bar{g}_k - \nabla f(x_k))\|\leq \sqrt{m_D}\cdot\tau/4)}\right] \\
& \leq \frac{1}{2} \mathbb{E}_k\left[\|D_k^{1/2}\nabla f(\bx_k)\|^2 \cdot \boldsymbol{1}_{(\|\bar{g}_k\| \le \tau,\|D_k^{1/2}(\bar{g}_k - \nabla f(x_k))\|\leq \sqrt{m_D}\cdot\tau/4)}\right] \\
& \leq \frac{1}{2} \mathbb{E}_k\left[\|D_k^{1/2}\nabla f(\bx_k)\|^2 \cdot \boldsymbol{1}_{(\|\bar{g}_k\| \le \tau)}\right]
\end{align*}
Plugging this into the above formula, we have
\begin{align*}
 & \mathbb{E}_k\left[\nabla f(x_k)^T D_k \bar{g}_k \cdot \boldsymbol{1}_{(\|\bar{g}_k\| \le \tau)}\right] \\
& \geq \frac{1}{2}\mathbb{E}_k\left[\| D_k^{1/2} \nabla f(x_k)\|^2\cdot \boldsymbol{1}_{(\|\bar{g}_k\| \le \tau)}\right] - \mathbb{E}_k\left[\|D_k^{1/2}\nabla f(x_k)\| \|D_k^{1/2} (\bar{g}_k - \nabla f(x_k))\| \cdot \boldsymbol{1}_{(\|D_k^{1/2}(\bar{g}_k - \nabla f(x_k))\|> \sqrt{m_D}\cdot\tau/4)}\right]\\
& \geq \frac{1}{2}\mathbb{E}_k\left[\| D_k^{1/2} \nabla f(x_k)\|^2\right]\cdot \mathbb{E}_k\left[\boldsymbol{1}_{(\|\bar{g}_k\| \le \tau)}\right] - \sqrt{M_D}\|\nabla f(x_k)\| \mathbb{E}_k[\|D_k^{1/2} (\bar{g}_k - \nabla f(x_k))\| \cdot \boldsymbol{1}_{(\|D_k^{1/2}(\bar{g}_k - \nabla f(x_k))\|> \sqrt{m_D}\cdot\tau/4)}]\\
& = \frac{p_g}{2}\mathbb{E}_k[\| D_k^{1/2} \nabla f(x_k)\|^2] - \sqrt{M_D}\|\nabla f(x_k)\| \mathbb{E}_k[\|D_k^{1/2} (\bar{g}_k - \nabla f(x_k))\| \cdot \boldsymbol{1}_{(\|D_k^{1/2}(\bar{g}_k - \nabla f(x_k))\|> \sqrt{m_D}\cdot\tau/4)}].
\end{align*}

Next, we bound $\mathbb{E}_k[\|D_k^{1/2} (\bar{g}_k - \nabla f(x_k))\| \cdot \boldsymbol{1}_{(\|D_k^{1/2}(\bar{g}_k - \nabla f(x_k))\|> \sqrt{m_D}\cdot\tau/4)}]$.

\begin{align*}
    M_D^{p/2} \sigma^p & \geq M_D^{p/2}\mathbb{E}_k[\|\barg_k-\nabla f(\bx_k)\|^p]\geq \mathbb{E}_k[\|D_k^{1/2}(\barg_k-\nabla f(\bx_k))\|^p] \\
    & = \mathbb{E}_k[\|D_k^{1/2}(\barg_k-\nabla f(\bx_k))\| \cdot \|D_k^{1/2}(\barg_k-\nabla f(\bx_k))\|^{p-1}] \\
    & \geq \mathbb{E}_k[\|D_k^{1/2}(\barg_k-\nabla f(\bx_k))\| \cdot \|D_k^{1/2}(\barg_k-\nabla f(\bx_k))\|^{p-1}\cdot \boldsymbol{1}_{(\|D_k^{1/2}(\bar{g}_k - \nabla f(x_k))\|> \sqrt{m_D}\cdot\tau/4)}] \\
    & \geq \mathbb{E}_k[\|D_k^{1/2}(\barg_k-\nabla f(\bx_k))\| \cdot \left(\sqrt{m_D}\cdot\frac{\tau}{4}\right)^{p-1}\cdot \boldsymbol{1}_{(\|D_k^{1/2}(\bar{g}_k - \nabla f(x_k))\|> \sqrt{m_D}\cdot\tau/4)}].
\end{align*}
Rearranging the terms yields
\begin{align*}
    \mathbb{E}_k[\|D_k^{1/2}(\barg_k-\nabla f(\bx_k))\|  \cdot \boldsymbol{1}_{(\|D_k^{1/2}(\bar{g}_k - \nabla f(x_k))\|> \sqrt{m_D}\cdot\tau/4)}] \leq \frac{M_D^{p/2} \sigma^p}{\left(\sqrt{m_D}\cdot\frac{\tau}{4}\right)^{p-1}}.
\end{align*}
Thus
\begin{align*}
    \mathbb{E}_k[\nabla f(x_k)^T D_k \bar{g}_k \cdot \boldsymbol{1}_{(\|\bar{g}_k\| \le \tau)}]
    \geq \frac{p_g}{2}\mathbb{E}_k[\| D_k^{1/2} \nabla f(x_k)\|^2] -  \|\nabla f(\bx_k)\|\frac{M_D^{p/2+1} \sigma^p}{\left(\sqrt{m_D}\cdot\frac{\tau}{4}\right)^{p-1}} .
\end{align*}

Next, we examine $\tau \cdot \mathbb{E}_k\left[\nabla f(x_k)^T D_k \frac{\bar{g}_k}{\|\bar{g}_k\|} \cdot \boldsymbol{1}_{(\|\bar{g}_k\| > \tau)} \right]$.

We first consider the case $\|D_k^{1/2}\nabla f(\bx_k)\|\geq 2\|D_k^{1/2}(\barg_k-\nabla f(\bx_k))\|$, where we have
\begin{align*}
    \nabla f(\bx_k)^TD_k\barg_k 
    &= \nabla f(\bx_k)^TD_k\nabla f(\bx_k) +\nabla f(\bx_k)^TD_k (\barg_k - \nabla f(\bx_k)) \\
    & \geq \nabla f(\bx_k)^TD_k\nabla f(\bx_k) - \|D_k^{1/2}\nabla f(\bx_k)\|\|D_k^{1/2} (\barg_k - \nabla f(\bx_k))\| \\
    & \geq \nabla f(\bx_k)^TD_k\nabla f(\bx_k) - \frac{1}{2}\|D_k^{1/2}\nabla f(\bx_k)\|^2 \\
    & = \frac{1}{2}\|D_k^{1/2}\nabla f(\bx_k)\|^2.
\end{align*}

In addition, 
\begin{align*}
    \sqrt{m_D}\|\nabla f(\bx_k) + \barg_k - \nabla f(\bx_k)\| & \leq \|D_k^{1/2}(\nabla f(\bx_k) + \barg_k - \nabla f(\bx_k))\|\\
    & \leq \|D_k^{1/2}\nabla f(\bx_k)\| + \| D_k^{1/2}(\barg_k - \nabla f(\bx_k))\| \\
    & \leq \frac{3}{2}\|D_k^{1/2}\nabla f(\bx_k)\|.
\end{align*}
Rearranging the terms yields
\begin{equation*}
   \|\nabla f(\bx_k) + \barg_k - \nabla f(\bx_k)\| \leq  \frac{3}{2\sqrt{m_D}}\|D_k^{1/2}\nabla f(\bx_k)\|.
\end{equation*}

Therefore, we have
\begin{align*}
    \nabla f(\bx_k)^TD_k\frac{\barg_k}{\|\barg_k\|} & = \frac{\nabla f(\bx_k)^TD_k(\nabla f(\bx_k)+\barg_k-\nabla f(\bx_k))}{\|\nabla f(\bx_k)+\barg_k-\nabla f(\bx_k)\|}\\
    & \geq \frac{\sqrt{m_D}}{3}\|D_k^{1/2}\nabla f(\bx_k)\|.
\end{align*}


When $\|D_k^{1/2}\nabla f(\bx_k)\|<2\|D_k^{1/2}(\barg_k-\nabla f(\bx_k))\|$, we have
\begin{align*}
    \nabla f(\bx_k)^TD_k\frac{\barg_k}{\|\barg_k\|} & \geq - \| D_k^{1/2}\nabla f(\bx_k)\| \frac{\| D_k^{1/2}\barg_k\|}{\|\barg_k\|}\\
   & \geq - \sqrt{M_D}\| D_k^{1/2}\nabla f(\bx_k)\|\\
   & =  \frac{\sqrt{m_D}}{3}\| D_k^{1/2}\nabla f(\bx_k)\| - \left( \frac{\sqrt{m_D}}{3} + \sqrt{M_D}\right)\| D_k^{1/2}\nabla f(\bx_k)\| \\
   & \geq \frac{\sqrt{m_D}}{3}\| D_k^{1/2}\nabla f(\bx_k)\| - 2\left( \frac{\sqrt{m_D}}{3} + \sqrt{M_D}\right)\| D_k^{1/2}(\barg_k-\nabla f(\bx_k))\|.
\end{align*}
Combining the above formulas, we have
\begin{align*}
   & \tau \cdot \mathbb{E}_k\left[\nabla f(x_k)^T D_k \frac{\bar{g}_k}{\|\bar{g}_k\|} \cdot \boldsymbol{1}_{(\|\bar{g}_k\| > \tau)} \right] \\
   & \geq  
   \tau \cdot \mathbb{E}_k\left[\left(\frac{\sqrt{m_D}}{3}\| D_k^{1/2}\nabla f(\bx_k)\| - 2\left( \frac{\sqrt{m_D}}{3} + \sqrt{M_D}\right)\| D_k^{1/2}(\barg_k-\nabla f(\bx_k))\|\right) \cdot \boldsymbol{1}_{(\|\bar{g}_k\| > \tau)} \right] \\
   & = \tau \frac{\sqrt{m_D}}{3} \cdot \mathbb{E}_k\left[\| D_k^{1/2}\nabla f(\bx_k)\| \cdot \boldsymbol{1}_{(\|\bar{g}_k\| > \tau)} \right]  \\
   & \quad - 2\left( \frac{\sqrt{m_D}}{3} + \sqrt{M_D}\right) \tau \cdot \mathbb{E}_k\left[ \| D_k^{1/2}(\barg_k-\nabla f(\bx_k))\|\cdot \boldsymbol{1}_{(\|\bar{g}_k\| > \tau) } \right] \\
   & \geq \tau \frac{\sqrt{m_D}}{3} \cdot \mathbb{E}_k\left[\| D_k^{1/2}\nabla f(\bx_k)\|\right] \cdot \mathbb{E}_k\left[\boldsymbol{1}_{(\|\bar{g}_k\| > \tau) }\right]  \\
   & \quad - 2\tau \left( \frac{\sqrt{m_D}}{3} + \sqrt{M_D}\right)  \cdot \mathbb{E}_k\left[ \| D_k^{1/2}(\barg_k-\nabla f(\bx_k))\|\cdot \boldsymbol{1}_{(\|\bar{g}_k\| > \tau)} \right] \\
   & \geq \tau (1-p_g) \frac{\sqrt{m_D}}{3} \cdot \mathbb{E}_k\left[\| D_k^{1/2}\nabla f(\bx_k)\|\right]  - 2\tau \left( \frac{\sqrt{m_D}}{3} + \sqrt{M_D}\right)  \cdot \mathbb{E}_k\left[ \| D_k^{1/2}(\barg_k-\nabla f(\bx_k))\|  \right] \\
   & \geq \tau (1-p_g) \frac{\sqrt{m_D}}{3} \cdot \mathbb{E}_k\left[\| D_k^{1/2}\nabla f(\bx_k)\|\right]  - 2\tau \left( \frac{\sqrt{m_D}}{3} + \sqrt{M_D}\right) \sqrt{M_D} \cdot \sigma .
\end{align*}

Combining the above derivations, we have
\begin{align*}
    \mathbb{E}_k[\nabla f(x_k)^T D_k \widehat{\barg}_k ] & = \mathbb{E}_k[\nabla f(x_k)^T D_k \bar{g}_k \cdot \boldsymbol{1}_{(\|\bar{g}_k\| \le \tau)} ]  + \tau \cdot \mathbb{E}_k\left[\nabla f(x_k)^T D_k \frac{\bar{g}_k}{\|\bar{g}_k\|} \cdot \boldsymbol{1}_{(\|\bar{g}_k\| > \tau)} \right] \\
    & \geq \frac{p_g}{2}\mathbb{E}_k[\| D_k^{1/2} \nabla f(x_k)\|^2] -  \|\nabla f(\bx_k)\|\frac{M_D^{p/2+1} \sigma^p}{\left(\sqrt{m_D}\cdot\frac{\tau}{4}\right)^{p-1}} \\
    & \quad + \tau (1-p_g) \frac{\sqrt{m_D}}{3} \cdot \mathbb{E}_k\left[\| D_k^{1/2}\nabla f(\bx_k)\|\right]  - 2\tau \left( \frac{\sqrt{m_D}}{3} + \sqrt{M_D}\right) \sqrt{M_D} \cdot \sigma .
\end{align*}
Note that since $\|D_k^{1/2}\nabla f(\bx_k)\|\geq \sqrt{m_D}\|\nabla f(\bx_k)\|$, we have
\begin{align*}
    \mathbb{E}_k[\nabla f(x_k)^T D_k \widehat{\barg}_k ]
    & \geq \frac{p_g}{2}\sqrt{m_D}\mathbb{E}_k[\| D_k^{1/2} \nabla f(x_k)\|\|\nabla f(\bx_k)\|] -  \|\nabla f(\bx_k)\|\frac{M_D^{p/2+1} \sigma^p}{\left(\sqrt{m_D}\cdot\frac{\tau}{4}\right)^{p-1}} \\
    & \quad + \tau (1-p_g) \frac{\sqrt{m_D}}{4} \cdot \mathbb{E}_k\left[\| D_k^{1/2}\nabla f(\bx_k)\|\right]  - 2\tau \left( \frac{\sqrt{m_D}}{3} + \sqrt{M_D}\right) \sqrt{M_D} \cdot \sigma .
\end{align*}
Since $\|\nabla f(\bx_k)\|\geq \tau/2$, we further have
\begin{align*}
    \mathbb{E}_k[\nabla f(x_k)^T D_k \widehat{\barg}_k ]
    & \geq \frac{p_g}{4}\sqrt{m_D}\tau\cdot\mathbb{E}_k[\| D_k^{1/2} \nabla f(x_k)\|] -  \|\nabla f(\bx_k)\|\frac{M_D^{p/2+1} \sigma^p}{\left(\sqrt{m_D}\cdot\frac{\tau}{4}\right)^{p-1}} \\
    & \quad + \tau (1-p_g) \frac{\sqrt{m_D}}{4} \cdot \mathbb{E}_k\left[\| D_k^{1/2}\nabla f(\bx_k)\|\right]  - 2\tau \left( \frac{\sqrt{m_D}}{3} + \sqrt{M_D}\right) \sqrt{M_D} \cdot \sigma \\
    & = \tau  \frac{\sqrt{m_D}}{4} \cdot \mathbb{E}_k\left[\| D_k^{1/2}\nabla f(\bx_k)\|\right] -  \|\nabla f(\bx_k)\|\frac{M_D^{p/2+1} \sigma^p}{\left(\sqrt{m_D}\cdot\frac{\tau}{4}\right)^{p-1}} - 2\tau \left( \frac{\sqrt{m_D}}{3} + \sqrt{M_D}\right) \sqrt{M_D} \cdot \sigma \\
    & \geq \tau  \frac{m_D}{4} \cdot \| \nabla f(\bx_k)\| -  \|\nabla f(\bx_k)\|\frac{M_D^{p/2+1} \sigma^p}{\left(\sqrt{m_D}\cdot\frac{\tau}{4}\right)^{p-1}} - 2\tau \left( \frac{\sqrt{m_D}}{3} + \sqrt{M_D}\right) \sqrt{M_D} \cdot \sigma.
\end{align*}
Since we set 
\begin{equation*}
    \tau \geq \max\left\{2,  4\cdot 3^{1/p}\cdot \sigma\left(\frac{M_D}{m_D}\right)^{\frac{p+1}{2p}}, 64 \sigma \frac{M_D}{m_D} \right\},
\end{equation*}
we further have
\begin{equation*}
    \frac{M_D^{p/2+1} \sigma^p}{\left(\sqrt{m_D}\cdot\frac{\tau}{4}\right)^{p-1}} \leq \frac{\tau m_D}{12}, \quad 2\tau \left( \frac{\sqrt{m_D}}{3} + \sqrt{M_D}\right) \sqrt{M_D} \cdot \sigma \leq \frac{\tau^2}{24}m_D.
\end{equation*}

These imply
\begin{align*}
    \mathbb{E}_k[\nabla f(x_k)^T D_k \widehat{\barg}_k ]
    & \geq   \frac{\tau m_D}{4} \cdot \| \nabla f(\bx_k)\| -  \frac{\tau m_D}{12}\|\nabla f(\bx_k)\| - \frac{\tau^2}{24}m_D\\
    & \geq \frac{\tau m_D}{4} \cdot \| \nabla f(\bx_k)\| -  \frac{\tau m_D}{6}\|\nabla f(\bx_k)\|\\
    & = \frac{\tau m_D}{12}\|\nabla f(\bx_k)\|.
\end{align*}

Therefore, 
 \begin{align*}
\mathbb{E}_k[f(\bx_{k+1})]  & \leq f(\bx_k) - \eta  \mathbb{E}_k[\nabla f(\bx_k)^T D_k\widehat{\barg}_k] + \frac{\eta^2 L }{2} \mathbb{E}_k[\|D_k\widehat{\barg}_k\|^2]\\
& \leq - \eta \frac{\tau m_D}{12}\|\nabla f(\bx_k)\| + \frac{\eta^2 L }{2}\tau^2 \\
& \leq - \eta \frac{\tau m_D}{12}\|\nabla f(\bx_k)\| + \eta^2 L  \tau \|\nabla f(\bx_k)\| \\
& \leq - \eta \frac{\tau m_D}{24}\|\nabla f(\bx_k)\| \\
& \leq - \eta \frac{ m_D}{12}\|\nabla f(\bx_k)\|.
\end{align*}
We thus complete the proof.
\end{proof}

\end{document}
